% problem_id: S001-lemma-smooth-test-artinian
% source_id: S001 (Stacks Project)
% source_file: algebra.tex
% statement_locator: algebra.tex lines 38949--38974
% proof_locator: algebra.tex lines 38976--39139
% env: lemma
% id_note: label 在本来源内唯一
% pairing: tight (verified)
% status: complete (statement + author proof verbatim + reason + locators)
%
\begin{lemma}
\label{lemma-smooth-test-artinian}
Let $R \to S$ be a ring map. Let $\mathfrak q$ be a prime ideal of
$S$ lying over $\mathfrak p \subset R$. Assume $R$ is Noetherian
and $R \to S$ of finite type.
The following are equivalent:
\begin{enumerate}
\item $R \to S$ is smooth at $\mathfrak q$,
\item for every surjection of local $R$-algebras
$(B', \mathfrak m') \to (B, \mathfrak m)$
with $\Ker(B' \to B)$ having square zero
and every solid commutative diagram
$$
\xymatrix{
S \ar[r] \ar@{-->}[rd] & B \\
R \ar[r] \ar[u] & B' \ar[u]
}
$$
such that $\mathfrak q = S \cap \mathfrak m$ there exists a dotted
arrow making the diagram commute,
\item same as in (2) but with $B' \to B$ ranging over small extensions, and
\item same as in (2) but with $B' \to B$ ranging over small extensions
such that in addition $S \to B$ induces an isomorphism
$\kappa(\mathfrak q) \cong \kappa(\mathfrak m)$.
\end{enumerate}
\end{lemma}

\begin{proof}
Assume (1). This means there exists a $g \in S$, $g \not \in \mathfrak q$
such that $R \to S_g$ is smooth. By
Proposition \ref{proposition-smooth-formally-smooth}
we know that $R \to S_g$ is formally smooth. Note that given any diagram
as in (2) the map $S \to B$ factors automatically through $S_{\mathfrak q}$
and a fortiori through $S_g$. The formal smoothness of $S_g$ over $R$
gives us a morphism $S_g \to B'$ fitting into a similar diagram with $S_g$ at
the upper left corner. Composing with $S \to S_g$ gives the desired arrow.
In other words, we have shown that (1) implies (2).

\medskip\noindent
Clearly (2) implies (3) and (3) implies (4).

\medskip\noindent
Assume (4). We are going to show that (1) holds, thereby finishing the
proof of the lemma. Choose a presentation
$S = R[x_1, \ldots, x_n]/(f_1, \ldots, f_m)$.
This is possible as $S$ is of finite type over $R$ and therefore of finite
presentation (see
Lemma \ref{lemma-Noetherian-finite-type-is-finite-presentation}).
Set $I = (f_1, \ldots, f_m)$.
Consider the naive cotangent complex
$$
\text{d} : I/I^2
\longrightarrow
\bigoplus\nolimits_{j = 1}^n S\text{d}x_j
$$
of this presentation (see Section \ref{section-netherlander}).
It suffices to show that when we localize this complex at $\mathfrak q$
then the map becomes a split injection, see Lemma \ref{lemma-smooth-at-point}.
Denote $S' = R[x_1, \ldots, x_n]/I^2$.
By Lemma \ref{lemma-differential-mod-power-ideal} we have
$$
S \otimes_{S'} \Omega_{S'/R} =
S \otimes_{R[x_1, \ldots, x_n]} \Omega_{R[x_1, \ldots, x_n]/R} =
\bigoplus\nolimits_{j = 1}^m S\text{d}x_j.
$$
Thus the map
$$
\text{d} :
I/I^2
\longrightarrow
S \otimes_{S'} \Omega_{S'/R}
$$
is the same as the map in the naive cotangent complex above. In particular
the truth of the assertion we are trying to prove
depends only on the three rings $R \to S' \to S$.
Let $\mathfrak q' \subset R[x_1, \ldots, x_n]$ be the prime ideal
corresponding to $\mathfrak q$. Since
localization commutes with taking modules of differentials
(Lemma \ref{lemma-differentials-localize}) we see that it suffices to show
that the map
\begin{equation}
\label{equation-target-map}
\text{d} :
I_{\mathfrak q'}/I_{\mathfrak q'}^2
\longrightarrow
S_{\mathfrak q} \otimes_{S'_{\mathfrak q'}} \Omega_{S'_{\mathfrak q'}/R}
\end{equation}
coming from $R \to S'_{\mathfrak q'} \to S_{\mathfrak q}$
is a split injection.

\medskip\noindent
Let $N \in \mathbf{N}$ be an integer.
Consider the ring
$$
B'_N = S'_{\mathfrak q'} / (\mathfrak q')^N S'_{\mathfrak q'}
= (S'/(\mathfrak q')^N S')_{\mathfrak q'}
$$
and its quotient $B_N = B'_N/IB'_N$. Note that
$B_N \cong S_{\mathfrak q}/\mathfrak q^NS_{\mathfrak q}$.
Observe that $B'_N$ is an Artinian local ring since it is the
quotient of a local Noetherian ring by a power of its maximal ideal.
Consider a filtration of the kernel $I_N$ of $B'_N \to B_N$
by $B'_N$-submodules
$$
0 \subset J_{N, 1} \subset J_{N, 2} \subset \ldots \subset J_{N, n(N)} = I_N
$$
such that each successive quotient $J_{N, i}/J_{N, i - 1}$ has length $1$.
(As $B'_N$ is Artinian such a filtration exists.)
This gives a sequence of small extensions
$$
B'_N  \to B'_N/J_{N, 1} \to B'_N/J_{N, 2} \to \ldots \to
B'_N/J_{N, n(N)} = B'_N/I_N
= B_N = S_{\mathfrak q}/\mathfrak q^NS_{\mathfrak q}
$$
Applying condition (4) successively to these small extensions
starting with the map $S \to B_N$ we see there
exists a commutative diagram
$$
\xymatrix{
S \ar[r] \ar[rd] & B_N  \\
R \ar[r] \ar[u] & B'_N \ar[u]
}
$$
Clearly the ring map $S \to B'_N$ factors as $S \to S_{\mathfrak q} \to B'_N$
where $S_{\mathfrak q} \to B'_N$ is a local homomorphism of local rings.
Moreover, since the maximal ideal of $B'_N$ to the $N$th power is zero
we conclude that $S_{\mathfrak q} \to B'_N$ factors through
$S_{\mathfrak q}/(\mathfrak q)^NS_{\mathfrak q} = B_N$. In other words
we have shown that for all $N \in \mathbf{N}$ the surjection of
$R$-algebras $B'_N \to B_N$ has a splitting.

\medskip\noindent
Consider the presentation
$$
I_N \to B_N \otimes_{B'_N} \Omega_{B'_N/R} \to \Omega_{B_N/R} \to 0
$$
coming from the surjection $B'_N \to B_N$ with kernel $I_N$ (see
Lemma \ref{lemma-differential-seq}). By the above the $R$-algebra map
$B'_N \to B_N$ has a right inverse. Hence by
Lemma \ref{lemma-differential-seq-split} we see that the sequence above
is split exact! Thus for every $N$ the map
$$
I_N \longrightarrow B_N \otimes_{B'_N} \Omega_{B'_N/R}
$$
is a split injection. The rest of the proof is gotten by unwinding what
this means exactly. Note that
$$
I_N = I_{\mathfrak q'}/
(I_{\mathfrak q'}^2 + (\mathfrak q')^N \cap I_{\mathfrak q'})
$$
By Artin-Rees (Lemma \ref{lemma-Artin-Rees}) we find a $c \geq 0$
such that
$$
S_{\mathfrak q}/\mathfrak q^{N - c}S_{\mathfrak q}
\otimes_{S_{\mathfrak q}} I_N =
S_{\mathfrak q}/\mathfrak q^{N - c}S_{\mathfrak q}
\otimes_{S_{\mathfrak q}}
I_{\mathfrak q'}/I_{\mathfrak q'}^2
$$
for all $N \geq c$
(these tensor product are just a fancy way of dividing by
$\mathfrak q^{N - c}$). We may of course assume $c \geq 1$.
By Lemma \ref{lemma-differential-mod-power-ideal} we see that
$$
S'_{\mathfrak q'}/(\mathfrak q')^{N - c}S'_{\mathfrak q'}
\otimes_{S'_{\mathfrak q'}} \Omega_{B'_N/R} =
S'_{\mathfrak q'}/(\mathfrak q')^{N - c}S'_{\mathfrak q'}
\otimes_{S'_{\mathfrak q'}} \Omega_{S'_{\mathfrak q'}/R}
$$
we can further tensor this by $B_N = S_{\mathfrak q}/\mathfrak q^N$
to see that
$$
S_{\mathfrak q}/\mathfrak q^{N - c}S_{\mathfrak q}
\otimes_{S'_{\mathfrak q'}} \Omega_{B'_N/R} =
S_{\mathfrak q}/\mathfrak q^{N - c}S_{\mathfrak q}
\otimes_{S'_{\mathfrak q'}} \Omega_{S'_{\mathfrak q'}/R}.
$$
Since a split injection remains a split injection after tensoring
with anything we see that
$$
S_{\mathfrak q}/\mathfrak q^{N - c}S_{\mathfrak q}
\otimes_{S_{\mathfrak q}}
(\ref{equation-target-map}) =
S_{\mathfrak q}/\mathfrak q^{N - c}S_{\mathfrak q}
\otimes_{S_{\mathfrak q}/\mathfrak q^N S_{\mathfrak q}}
(I_N \longrightarrow B_N \otimes_{B'_N} \Omega_{B'_N/R})
$$
is a split injection for all $N \geq c$. By
Lemma \ref{lemma-split-injection-after-completion} we see that
(\ref{equation-target-map}) is a split injection. This finishes the proof.
\end{proof}
